\chapter{决策摘要：全量重筛后的机会曲线}

截至2026年7月15日，官方分页半年报预告表共4429条指标，排除B股后保留4339条A股指标，透视为1680家公司，其中1678家完成申万行业映射。统一规则筛出142只高影响候选，按价格位置、扣非纯度、Q1现金流和研报时效形成39只优先池。

\begin{houseviewbox}[AStock House View]
本次刷新不是把旧名单延长，而是重新从官方全量表开始。正式当前价模型自动筛出3只：恒逸石化（000703）、宏桥控股（002379）、川能动力（000155）。它们只代表刷新后的模型候选，不等同于不经IC复核的买入建议。
\end{houseviewbox}

\section{正式当前价模型}

\begin{exhibitbox}[Exhibit 1：2026-07-15正式模型候选]
\scriptsize
\begin{tabularx}{\textwidth}{L{1.0cm}L{1.6cm}L{1.6cm}R{0.9cm}R{0.9cm}R{0.9cm}R{0.9cm}R{0.85cm}X}
\toprule
\textbf{代码} & \textbf{公司} & \textbf{行业} & \textbf{现价} & \textbf{熊} & \textbf{基准} & \textbf{概率值} & \textbf{空间} & \textbf{动作} \\
\midrule
000703 & 恒逸石化 & 石油石化 & 14.48 & 10.86 & 25.59 & 23.42 & 61.7\% & 高空间验证 \\
002379 & 宏桥控股 & 有色金属 & 18.66 & 14.00 & 26.24 & 25.80 & 38.3\% & 高空间验证 \\
000155 & 川能动力 & 公用事业 & 11.09 & 5.43 & 9.03 & 10.04 & -9.5\% & 估值观察 \\
\bottomrule
\end{tabularx}
\normalsize
\end{exhibitbox}

\section{全量筛选结构}

\noindent
\begin{tabularx}{\textwidth}{L{4.0cm}R{1.5cm}X}
\toprule
\textbf{层级} & \textbf{数量} & \textbf{解释} \\
\midrule
官方预告公司 & 1680 & 2026H1归母预告主键，扣非/EPS/营收为质量字段 \\
高影响候选 & 142 & 利润规模、同比、扣非纯度与行业阶段统一筛选 \\
优先池 & 39 & 当前价格位置、Q1财务、现金流、研报时效二次筛选 \\
可定价候选 & 141 & 业务匹配的熊/基准/牛筛选区间 \\
正式模型候选 & 3 & 当前价格、股本、情景、质量和动作均已结构化 \\
\bottomrule
\end{tabularx}

候选处置分布为：静默吸纳优先 1只、低位盈利优先 34只、启动后仍有空间 4只、盈利验证观察 83只、业绩兑现、价格先行 11只、历史不足观察 1只、预减观察 2只、非经常主导排除 6只。行业阶段标签中，单日反弹 1063家、启动确认 50家、资金观察 165家、低位无流 147家、静默吸纳 128家、无信号 125家、未映射 2家；行业日报/周报均为31/31完整表，但观察日分别为7月14日和截至7月10日；连续流入公开表仅30/112，因此不把它扩展成完整112只名单。
